\renewcommand{\backrefpagesname}{Citado na(s) página(s):~}
\renewcommand{\backref}{}
\renewcommand*{\backrefalt}[4]{
	\ifcase #1 %
		Nenhuma citação no texto.%
	\or
		Citado na página #2.%
	\else
		Citado #1 vezes nas páginas #2.%
	\fi}%
% ---

\usetikzlibrary{shapes,snakes}
\usetikzlibrary{calc}

\theoremstyle{definition}
\newtheorem{myex}{Exemplo}
\newtheorem{mydef}{Definição}

\providecommand{\tightlist}{%
	  \setlength{\itemsep}{0pt}\setlength{\parskip}{0pt}}

\lstset{%
language=C++,           %linguagem
numbers=left,           %posição dos números
stepnumber=1,           %frequencia de aparição dos números
numbersep=5pt,
%numberstyle=\zebra{gray!15}{white!35},
%basewidth={0.6em,0.45em},
%fontadjust=true,
%mathescape=true,
tabsize=4,
commentstyle=\color{blue},
literate={á}{{\'a}}1 {à}{{\`a}}1 {ã}{{\~a}}1 {é}{{\'e}}1 {É}{{\'E}}1 {ê}{{\^e}}1 {õ}{{\~o}}1 {í}{{\'i}}1 {ó}{{\'o}}1 {ú}{{\'u}}1 {ç}{{\c c}}1 {³}{{$^3$}}1 {Ω}{{$\Omega$}}1
%breaklines=true,
%showstringspaces=false,
%stringstyle=\color{cyan},
%basicstyle=\small\ttfamily
}
\newcommand{\curso}[1]{\def\imprimircurso{#1}}

\newcommand{\palavraChaveUm}[1]{\def\imprimirpalavrachaveum{#1}}
\newcommand{\palavraChaveDois}[1]{\def\imprimirpalavrachavedois{#1}}

\newcommand{\cdu}[1]{\def\nomecdu{#1}}
\newcommand{\dataDaAprovacao}[1]{\def\imprimirdatadaaprovacao{#1}}

\newcommand{\membroConvidadoUm}[1]{\def\imprimirmembroconvidadoum{#1}}
\newcommand{\membroConvidadoDois}[1]{\def\imprimirmembroconvidadodois{#1}}

\newcommand\BackgroundPic{%
	\put(0,0){%
		\parbox[b][\paperheight]{\paperwidth}{%
			\vfill
			\centering
			\includegraphics[width=\paperwidth,height=\paperheight,%
				keepaspectratio]{img/capa.jpg}%
			\vfill
		}
	}
}

\renewcommand{\imprimircapa}{%
  \begin{capa}%
    \center
	\AddToShipoutPicture*{\BackgroundPic}

%	\begin{huge}
%		\textbf{\textsc{Trabalho de Conclusão de Curso}}
%	\end{huge}

    \vspace*{2.7in}
	{\textbf{\large\imprimirinstituicao}}
	\par
	{\textbf{\large\imprimircurso}}

	\vspace{0.5in}

    {\ABNTEXchapterfont\bfseries\LARGE\imprimirtitulo}
    \vspace*{\fill}

	\begin{flushright}
    	\textbf{{\large{Autor: \imprimirautor}}}
		\par
    	\textbf{{\large{Orientador: \imprimirorientador}}}
		\par
		\textbf{{\large{Coorientador: \imprimircoorientador}}}

	\end{flushright}

    \vspace*{0.2in}
    \textbf{{\large\imprimirlocal}}
    \par
    \textbf{{\large\imprimirdata}}

    \vspace*{2.2in}
  \end{capa}
}
